\section{从随机游走到布朗运动}
\label{sec:06-Random-Processes/06-Diffusion-and-SDE/01}

把一天的股价按分钟画成折线，得到一条毛糙的曲线。嫌它糙，换成按秒采样重画——曲线没有变光滑，只是把同样的毛糙搬进了更小的尺度。再细到毫秒，还是那副样子。

这跟日常经验拧着。地球是圆的而脚下的路是平的；可微函数在足够小的邻域里就是它的切线；放大总该把曲线拉直。训练日志里也有同一件事：每 100 步记一次 loss，曲线抖；改成每 10 步记一次，抖动幅度小了，形状却没变。采样加密并没有让你看清趋势的细节，你看清的是噪声的细节。

\hyperref[sec:06-Random-Processes/04-Applications/01]{``随机游走：从局部步进到整体尺度''}处理过一个相邻的问题：无偏随机游走走 \(n\) 步之后，典型距离按 \(\sqrt n\) 增长。那里问的是终点落在哪。本节问的是整条轨迹：把步子迈得越来越密、越来越小，轨迹本身趋向什么？极限有名字，叫布朗运动（Brownian motion）。它是本单元后面每一个方程的地基，也是扩散生成模型里那个``噪声''的严格身份。

\subsection{只有一种缩放不会把极限毁掉}

设时间步长为 \(\Delta t\)，每步等概率移动 \(\pm\Delta x\)，各步独立。到时刻 \(t\) 共走 \(n=t/\Delta t\) 步（为省事设 \(t\) 是 \(\Delta t\) 的整数倍），位置为

\[
X^{(\Delta t)}_t=\Delta x\sum_{k=1}^{n}Z_k,
\]

其中 \(Z_k\) 取 \(\pm1\)，等概率且相互独立。均值为零，方差为

\[
\Var\big(X^{(\Delta t)}_t\big)=n(\Delta x)^2=\frac{(\Delta x)^2}{\Delta t}\,t.
\]

整个极限的命运压在比值 \((\Delta x)^2/\Delta t\) 上。若 \(\Delta x\) 比 \(\sqrt{\Delta t}\) 缩得快（比如取 \(\Delta x=\Delta t\)），比值趋于零，方差蒸发，极限是一个钉死在原点的过程；若缩得慢（比如 \(\Delta x\) 固定不动只缩时间），比值发散，粒子在任意短的时间里跑到无穷远。中间只剩一条缝：\(\Delta x=\sigma\sqrt{\Delta t}\)，方差稳定在 \(\sigma^2 t\)。

这个指数不是为了好看挑出来的。位移按步数的平方根增长，而缩小时间步会让步数按 \(1/\Delta t\) 增加，两者要互相抵消，单步位移就只能按 \(\sqrt{\Delta t}\) 缩。以下取 \(\sigma=1\)，称标准缩放。

\begin{center}
\begin{tikzpicture}[x=1cm,y=1cm,font=\small,>=stealth]
\begin{scope}
  \draw[black!45] (0,-0.95) -- (0,1.55);
  \draw[black!45] (0,0) -- (3.15,0);
  \draw[blue!75,thick] (0.00,0.00) -- (0.50,0.20) -- (1.00,0.61) -- (1.50,1.22) -- (2.00,1.02) -- (2.50,0.20) --
    (3.00,-0.61);
  \node[font=\footnotesize,align=center] at (1.5,-1.5) {\(\Delta t=1/6\)\\6 步};
\end{scope}
\begin{scope}[shift={(3.9,0)}]
  \draw[black!45] (0,-0.95) -- (0,1.55);
  \draw[black!45] (0,0) -- (3.15,0);
  \draw[black!22] (0.00,0.00) -- (0.50,0.20) -- (1.00,0.61) -- (1.50,1.22) -- (2.00,1.02) -- (2.50,0.20) --
    (3.00,-0.61);
  \draw[blue!75,thick] (0.00,0.00) -- (0.12,0.00) -- (0.25,0.20) -- (0.38,0.00) -- (0.50,0.20) -- (0.62,0.20) --
    (0.75,0.20) -- (0.88,0.41) -- (1.00,0.61) -- (1.12,0.82) -- (1.25,0.82) -- (1.38,0.82) --
    (1.50,1.22) -- (1.62,1.22) -- (1.75,1.22) -- (1.88,1.22) -- (2.00,1.02) -- (2.12,0.61) --
    (2.25,0.41) -- (2.38,0.20) -- (2.50,0.20) -- (2.62,-0.20) -- (2.75,-0.41) -- (2.88,-0.61) --
    (3.00,-0.61);
  \node[font=\footnotesize,align=center] at (1.5,-1.5) {\(\Delta t=1/24\)\\24 步};
\end{scope}
\begin{scope}[shift={(7.8,0)}]
  \draw[black!45] (0,-0.95) -- (0,1.55);
  \draw[black!45] (0,0) -- (3.15,0);
  \draw[black!22] (0.00,0.00) -- (0.50,0.20) -- (1.00,0.61) -- (1.50,1.22) -- (2.00,1.02) -- (2.50,0.20) --
    (3.00,-0.61);
  \draw[blue!75,thick] (0.00,0.00) -- (0.03,-0.10) -- (0.06,0.00) -- (0.09,0.10) -- (0.12,0.00) -- (0.16,0.10) --
    (0.19,0.20) -- (0.22,0.31) -- (0.25,0.20) -- (0.28,0.10) -- (0.31,0.00) -- (0.34,0.10) --
    (0.38,0.00) -- (0.41,0.10) -- (0.44,0.00) -- (0.47,0.10) -- (0.50,0.20) -- (0.53,0.31) --
    (0.56,0.20) -- (0.59,0.31) -- (0.62,0.20) -- (0.66,0.10) -- (0.69,0.20) -- (0.72,0.31) --
    (0.75,0.20) -- (0.78,0.31) -- (0.81,0.41) -- (0.84,0.51) -- (0.88,0.41) -- (0.91,0.51) --
    (0.94,0.61) -- (0.97,0.51) -- (1.00,0.61) -- (1.03,0.71) -- (1.06,0.82) -- (1.09,0.92) --
    (1.12,0.82) -- (1.16,0.92) -- (1.19,0.82) -- (1.22,0.92) -- (1.25,0.82) -- (1.28,0.92) --
    (1.31,0.82) -- (1.34,0.92) -- (1.38,0.82) -- (1.41,0.92) -- (1.44,1.02) -- (1.47,1.12) --
    (1.50,1.22) -- (1.53,1.12) -- (1.56,1.22) -- (1.59,1.12) -- (1.62,1.22) -- (1.66,1.12) --
    (1.69,1.02) -- (1.72,1.12) -- (1.75,1.22) -- (1.78,1.12) -- (1.81,1.22) -- (1.84,1.12) --
    (1.88,1.22) -- (1.91,1.12) -- (1.94,1.02) -- (1.97,0.92) -- (2.00,1.02) -- (2.03,0.92) --
    (2.06,0.82) -- (2.09,0.71) -- (2.12,0.61) -- (2.16,0.51) -- (2.19,0.41) -- (2.22,0.31) --
    (2.25,0.41) -- (2.28,0.51) -- (2.31,0.41) -- (2.34,0.31) -- (2.38,0.20) -- (2.41,0.10) --
    (2.44,0.20) -- (2.47,0.31) -- (2.50,0.20) -- (2.53,0.10) -- (2.56,0.00) -- (2.59,-0.10) --
    (2.62,-0.20) -- (2.66,-0.31) -- (2.69,-0.20) -- (2.72,-0.31) -- (2.75,-0.41) -- (2.78,-0.51) --
    (2.81,-0.61) -- (2.84,-0.51) -- (2.88,-0.61) -- (2.91,-0.71) -- (2.94,-0.82) -- (2.97,-0.71) --
    (3.00,-0.61);
  \node[font=\footnotesize,align=center] at (1.5,-1.5) {\(\Delta t=1/96\)\\96 步};
\end{scope}
\end{tikzpicture}

\emph{同一次实验记录得越来越密。步数翻四倍、步长减半（\(\Delta x=\sqrt{\Delta t}\)），灰线是最粗的那条骨架，蓝线是加密后的实际轨迹，两者始终穿过同样的节点。要点在纵向：幅度既没有塌成一条直线，也没有炸开，稳定在 \(\pm1\) 上下。这就是 \(\sqrt{\Delta t}\) 换来的东西。}
\end{center}

这套尺度关系在优化里照样成立。随机梯度下降每步走 \(\theta_{k+1}=\theta_k-\eta\nabla L(\theta_k)-\eta\,\xi_k\)，其中 \(\xi_k\) 是小批量带来的梯度噪声。把学习率 \(\eta\) 当作时间步，跑 \(K\) 步相当于走了 \(t=\eta K\) 的时间：确定性部分累积 \(\eta K=t\)，噪声部分累积的方差是 \(\eta^2K=\eta\,t\)。噪声的标准差因此正比于 \(\sqrt{\eta}\)——学习率同时扮演步长与噪声强度两个角色，改学习率不只是走快走慢，还在调节随机性的温度。下一节之后这条对应会写成方程。

\subsection{极限过程的四条性质}

固定时刻 \(t\)，步数 \(n=t/\Delta t\to\infty\)，中心极限定理直接给出

\[
X^{(\Delta t)}_t=\sqrt{\Delta t}\sum_{k=1}^{n}Z_k
\xrightarrow{\;d\;}\mathcal N(0,t).
\]

方差是 \(t\) 不是 \(1\)：时间越长，散布越宽。同理，区间 \([s,t]\) 上的位移由 \((t-s)/\Delta t\) 个小步合成，趋向 \(\mathcal N(0,t-s)\)；不相交的区间用的是不同的小步，因而相互独立。把折线在每个 \(\Delta t\) 内拉直，跳幅只有 \(\sqrt{\Delta t}\)，折角在极限中抹平，路径应当连续。

比中心极限定理更强的结论叫 Donsker 不变原理：只要单步零均值、方差有限，无论它是 \(\pm1\)、均匀分布还是别的什么，标准缩放后\emph{整条折线}的分布都收敛到同一个极限过程。本书只声明这一结论，不作证明——它要在函数空间上谈分布收敛。它值钱的地方是普适性：微观步进的形状被 \(\sqrt{\Delta t}\) 尺度洗得干干净净，极限只有一个，所以配得上一个专门的名字。

\begin{definition}
实值过程 \(W=(W_t)_{t\ge0}\) 称为标准布朗运动，若 \(W_0=0\)；不相交时间区间上的增量相互独立；对 \(s<t\) 有 \(W_t-W_s\sim\mathcal N(0,t-s)\)；且以概率一，\(t\mapsto W_t\) 连续。
\end{definition}

前三条是缩放极限的直接读法，第四条来自插值折线。第四条不能省：前三条只约束任意有限个时刻的联合分布，人为构造一个在随机时刻抖一下的坏版本，有限维分布与布朗运动完全一致，路径却到处跳。Kolmogorov 连续性定理保证满足前三条的过程总有一个连续版本（声明，不证），定义第四条挑的就是这个好版本。

由独立正态增量立刻得到 \(\Cov(W_s,W_t)=\min(s,t)\)，任意有限个时刻联合 Gaussian，所以布朗运动是 Gaussian 过程；未来增量与过去独立，所以它也是 Markov 过程。带初值、漂移与强度的版本写作 \(x+\mu t+\sigma W_t\)：\(\mu t\) 叠一个确定性趋势，\(\sigma\) 调节波动幅度。

\subsection{放大十倍，仍然同样粗糙}

回到开头那条毫秒级的股价曲线。布朗运动把这种``放不平''的性质推到了极端。

\begin{center}
\begin{tikzpicture}[x=1cm,y=1cm,font=\small,>=stealth]
\begin{scope}
  \draw[black!45] (0,-1.0) -- (0,1.05);
  \draw[blue!70] (0.00,0.78) -- (0.05,0.62) -- (0.10,0.75) -- (0.15,0.66) -- (0.20,0.59) -- (0.25,0.47) --
    (0.30,0.44) -- (0.35,0.41) -- (0.40,0.25) -- (0.45,0.44) -- (0.50,0.53) -- (0.55,0.50) --
    (0.60,0.44) -- (0.65,0.56) -- (0.70,0.56) -- (0.75,0.50) -- (0.80,0.34) -- (0.85,0.34) --
    (0.90,0.34) -- (0.95,0.16) -- (1.00,-0.06) -- (1.05,-0.03) -- (1.10,-0.03) -- (1.15,0.06) --
    (1.20,0.09) -- (1.25,0.28) -- (1.30,0.03) -- (1.35,0.25) -- (1.40,0.22) -- (1.45,0.16) --
    (1.50,0.22) -- (1.55,-0.03) -- (1.60,-0.09) -- (1.65,-0.06) -- (1.70,0.19) -- (1.75,0.25) --
    (1.80,0.06) -- (1.85,0.09) -- (1.90,0.03) -- (1.95,-0.03) -- (2.00,0.00) -- (2.05,0.00) --
    (2.10,-0.16) -- (2.15,-0.13) -- (2.20,-0.03) -- (2.25,-0.16) -- (2.30,-0.22) -- (2.35,-0.34) --
    (2.40,-0.25) -- (2.45,-0.09) -- (2.50,0.12) -- (2.55,0.19) -- (2.60,-0.09) -- (2.65,-0.13) --
    (2.70,-0.47) -- (2.75,-0.50) -- (2.80,-0.31) -- (2.85,-0.13) -- (2.90,0.03) -- (2.95,-0.03) --
    (3.00,0.03) -- (3.05,-0.25) -- (3.10,-0.59) -- (3.15,-0.63) -- (3.20,-0.78);
  \draw[red!70] (1.52,-0.30) rectangle (1.73,0.24);
  \node[font=\footnotesize] at (1.6,-1.35) {时间跨度 \(1\)};
\end{scope}
\draw[->,red!70] (3.55,0.0) -- (4.35,0.0);
\node[above,font=\footnotesize,red!70] at (3.95,0.05) {放大};
\begin{scope}[shift={(4.7,0)}]
  \draw[red!70] (0,-0.95) rectangle (3.2,0.95);
  \draw[blue!70] (0.00,0.06) -- (0.05,-0.06) -- (0.10,-0.19) -- (0.15,-0.06) -- (0.20,0.06) -- (0.25,-0.06) --
    (0.30,-0.19) -- (0.35,-0.06) -- (0.40,-0.19) -- (0.45,-0.31) -- (0.50,-0.19) -- (0.55,-0.31) --
    (0.60,-0.44) -- (0.65,-0.31) -- (0.70,-0.44) -- (0.75,-0.31) -- (0.80,-0.19) -- (0.85,-0.31) --
    (0.90,-0.44) -- (0.95,-0.56) -- (1.00,-0.44) -- (1.05,-0.56) -- (1.10,-0.44) -- (1.15,-0.56) --
    (1.20,-0.44) -- (1.25,-0.56) -- (1.30,-0.69) -- (1.35,-0.56) -- (1.40,-0.44) -- (1.45,-0.31) --
    (1.50,-0.19) -- (1.55,-0.31) -- (1.60,-0.44) -- (1.65,-0.31) -- (1.70,-0.19) -- (1.75,-0.31) --
    (1.80,-0.19) -- (1.85,-0.31) -- (1.90,-0.19) -- (1.95,-0.06) -- (2.00,0.06) -- (2.05,0.19) --
    (2.10,0.31) -- (2.15,0.44) -- (2.20,0.31) -- (2.25,0.19) -- (2.30,0.31) -- (2.35,0.19) --
    (2.40,0.06) -- (2.45,0.19) -- (2.50,0.31) -- (2.55,0.19) -- (2.60,0.31) -- (2.65,0.19) --
    (2.70,0.06) -- (2.75,0.19) -- (2.80,0.31) -- (2.85,0.44) -- (2.90,0.56) -- (2.95,0.69) --
    (3.00,0.56) -- (3.05,0.69) -- (3.10,0.56) -- (3.15,0.44) -- (3.20,0.31);
  \node[font=\footnotesize] at (1.6,-1.35) {时间跨度 \(1/64\)，纵向拉伸 \(8\) 倍};
\end{scope}
\end{tikzpicture}

\emph{右图是左图红框里那一小段的特写。时间轴放大 64 倍，纵轴只需放大 \(\sqrt{64}=8\) 倍，两张图就统计上无法区分。这正是 \(\Delta x\sim\sqrt{\Delta t}\) 的几何面孔：横竖两个方向按不同的幂次缩放，粗糙度才守恒。若纵轴也放大 64 倍，右图会撑出画面；若只放大 \(1.5\) 倍，右图会塌成一条直线。}
\end{center}

割线斜率的量级由此写死：

\[
\frac{W_{t+\Delta t}-W_t}{\Delta t}\ \sim\ \frac{\sqrt{\Delta t}}{\Delta t}=\frac{1}{\sqrt{\Delta t}}\longrightarrow\infty,
\]

而且符号随机翻转。\(dW/dt\) 不存在，``噪声的变化率''是一句没有指称的话。想写噪声驱动的方程 \(dX=\mu\,dt+\sigma\,dW\)，就得先把积分和链式法则重建一遍。

\subsection{位移的平方不肯消失}

不可微这件事可以量化得更漂亮。把 \([0,t]\) 分割成 \(0=t_0<\cdots<t_n=t\)，记 \(\Delta W_k=W_{t_{k+1}}-W_{t_k}\)、\(\Delta t_k=t_{k+1}-t_k\)，考察

\[
Q^{(n)}=\sum_{k=0}^{n-1}(\Delta W_k)^2 .
\]

期望是现成的：\(\E[(\Delta W_k)^2]=\Delta t_k\)，逐项加起来恰是 \(t\)。方差要用 Gaussian 的四阶矩，对 \(Z\sim\mathcal N(0,v)\) 有 \(\Var(Z^2)=3v^2-v^2=2v^2\)，再由增量独立，

\[
\Var\big(Q^{(n)}\big)=2\sum_k(\Delta t_k)^2\le 2\Big(\max_k\Delta t_k\Big)\sum_k\Delta t_k
=2t\max_k\Delta t_k\longrightarrow 0 .
\]

分割一细，随机性就被榨干了：\(Q^{(n)}\) 在均方意义下收敛到常数 \(t\)。这个极限记作 \([W]_t=t\)，叫二次变差。

\begin{center}
\begin{tikzpicture}[x=1cm,y=1cm,font=\small,>=stealth]
  \draw[->,black!55] (0,0) -- (3.75,0) node[right,font=\footnotesize] {\(t\)};
  \draw[->,black!55] (0,0) -- (0,3.35);
  \draw[black!45,dashed] (0,0) -- (3.4,2.2);
  \node[right,font=\footnotesize,black!60] at (3.42,2.2) {\(y=t\)};
  \draw[black!30] (0.00,0.00) -- (0.21,0.08) -- (0.42,0.34) -- (0.64,0.41) -- (0.85,0.43) -- (1.06,0.80) --
    (1.27,0.85) -- (1.49,0.89) -- (1.70,1.10) -- (1.91,1.15) -- (2.12,1.16) -- (2.34,1.16) --
    (2.55,1.27) -- (2.76,1.32) -- (2.98,1.43) -- (3.19,1.69) -- (3.40,3.14);
  \node[right,font=\footnotesize,black!45] at (3.42,3.10) {\(n=16\)};
  \draw[blue!40] (0.00,0.00) -- (0.11,0.09) -- (0.21,0.12) -- (0.32,0.15) -- (0.42,0.21) -- (0.53,0.31) --
    (0.64,0.32) -- (0.74,0.35) -- (0.85,0.41) -- (0.96,0.41) -- (1.06,0.60) -- (1.17,0.60) --
    (1.27,0.62) -- (1.38,0.83) -- (1.49,0.94) -- (1.59,0.96) -- (1.70,1.10) -- (1.81,1.24) --
    (1.91,1.33) -- (2.02,1.34) -- (2.12,1.35) -- (2.23,1.41) -- (2.34,1.43) -- (2.44,1.47) --
    (2.55,1.52) -- (2.66,1.68) -- (2.76,1.86) -- (2.87,2.13) -- (2.98,2.21) -- (3.08,2.34) --
    (3.19,2.35) -- (3.29,2.79) -- (3.40,2.84);
  \node[right,font=\footnotesize,blue!50] at (3.42,2.80) {\(n=64\)};
  \draw[blue!85,thick] (0.00,0.00) -- (0.07,0.06) -- (0.13,0.07) -- (0.20,0.11) -- (0.27,0.13) -- (0.33,0.15) --
    (0.40,0.18) -- (0.46,0.20) -- (0.53,0.24) -- (0.60,0.30) -- (0.66,0.33) -- (0.73,0.37) --
    (0.80,0.42) -- (0.86,0.46) -- (0.93,0.50) -- (1.00,0.55) -- (1.06,0.74) -- (1.13,0.76) --
    (1.20,0.78) -- (1.26,0.84) -- (1.33,0.92) -- (1.39,1.00) -- (1.46,1.05) -- (1.53,1.10) --
    (1.59,1.11) -- (1.66,1.19) -- (1.73,1.21) -- (1.79,1.30) -- (1.86,1.34) -- (1.93,1.39) --
    (1.99,1.42) -- (2.06,1.45) -- (2.12,1.48) -- (2.19,1.50) -- (2.26,1.53) -- (2.32,1.58) --
    (2.39,1.62) -- (2.46,1.68) -- (2.52,1.77) -- (2.59,1.83) -- (2.66,1.87) -- (2.72,1.93) --
    (2.79,1.97) -- (2.86,2.02) -- (2.92,2.09) -- (2.99,2.15) -- (3.05,2.22) -- (3.12,2.24) --
    (3.19,2.26) -- (3.25,2.36) -- (3.32,2.42) -- (3.39,2.49);
  \node[right,font=\footnotesize,blue!85] at (3.42,2.49) {\(n=256\)};
  \draw[red!65,thick] (0,0.05) -- (3.4,0.09);
  \node[above,font=\footnotesize,red!70] at (2.90,0.12) {光滑函数};
  \node[below,font=\footnotesize,align=center] at (1.7,-0.35) {累积平方位移 \(\sum_{k\le\,\cdot}(\Delta W_k)^2\)};
\end{tikzpicture}

\emph{同一条布朗路径，用三种粗细的分割累加平方位移。终值分别是 \(1.43\)、\(1.29\)、\(1.14\)，正朝虚线上的 \(t=1\) 靠拢，而且曲线越来越平顺——随机性随分割变细而退场。红线是一条光滑函数的同一个统计量，它趴在零附近不动。这张图上的两条走势，就是后面所有``反常''公式的分水岭。}
\end{center}

差别只在位移量级。光滑函数在小区间上的增量约为 \(f'(t)\Delta t\)，平方后是 \((\Delta t)^2\)，把 \(1/\Delta t\) 项加起来仍趋于零；布朗位移是 \(\sqrt{\Delta t}\)，平方是 \(\Delta t\)，同样多的项加起来恰好收敛到 \(t\)。

同一把尺子还判了另一件事的死刑。绝对值之和 \(\sum_k\abs{\Delta W_k}\) 每项量级 \(\sqrt{\Delta t}\)，共 \(t/\Delta t\) 项，合计 \(t/\sqrt{\Delta t}\to\infty\)。图中那条路径实测下来：分成 16 段时总变差是 \(3.8\)，64 段时 \(7.2\)，256 段时 \(13.4\)，1024 段时 \(23.8\)——每细分四倍就大约翻一倍，正是 \(\sqrt{\;}\) 的签名。布朗路径不是有界变差的，它的``弧长''在任何区间上都发散。而 Riemann--Stieltjes 积分 \(\int f\,dg\) 恰好要求积分器有界变差。于是 \(\int\cdot\,dW\) 不能逐条路径地按普通分割极限定义——不是精度不够，是定义本身不收敛。

\subsection{一阶随机，二阶确定}

把两件事并排放：\(W_t\) 本身是随机变量，均值零方差 \(t\)；\([W]_t\) 却是确定的数 \(t\)，一点随机性都不剩。一阶随机，二阶确定，这是布朗运动的代数指纹，也是本单元最该带走的一句话。

\begin{quote}
布朗位移与 \(\sqrt{\Delta t}\) 同阶，所以它的平方与 \(\Delta t\) 同阶。Taylor 展开里那个通常被当作高阶无穷小丢掉的二阶项，在这里和一阶项一样大——微积分的链式法则因此必须改写。
\end{quote}

这条因果链会串起后面四节。\(\sqrt{\Delta t}\) 尺度一头通向不可微与总变差发散，逼出下一节和之后的重建工作；另一头写成微分记号就是 \((dW)^2=dt\)，直接生出 Itô 公式里那个多出来的 \(\tfrac12 f''\sigma^2\)。读到那个修正项时若觉得突兀，回到这一节的最后两张图：它不是补丁，是 \(\sqrt{\Delta t}\) 的账单。

最后交代边界。布朗运动假设增量独立、正态、路径连续，三条在真实系统里都可能崩。金融市场有跳跃（闪崩）、有波动率聚集（大波动后跟着大波动，增量并不独立）、有重尾（极端收益远超正态预测）；物理扩散在分子碰撞的时间尺度上带记忆。它是连续世界的基准模型，跟随机游走在离散世界的地位一样——出发点，不是终点。

下一节先绕开单条路径的麻烦：不问某个粒子走到哪，只问一大群粒子形成的密度怎样演化。那里会出现一个让人踏实的反差——路径粗糙不堪，密度却服从一个干净的偏微分方程。
